\chapter{Design}

\section{Design Goals}

\textit{What should go here:} restate the specific teaching problem the design responds to. Connect the chapter back to the methodology by showing that the design exists to address the stated learning outcomes and hypotheses rather than to showcase VR for its own sake.

\section{Learning Outcomes as Design Drivers}

\textit{What should go here:} show how each learning outcome shaped the structure of the environment. This is where the dissertation should make explicit that the prototype was scoped around supported instruction execution tasks rather than around the whole of computer architecture.

\section{Pedagogical Approach}

\textit{What should go here:} explain the teaching model behind the system. This should cover guided tracing, progressive reduction of support, the role of embodied interaction, and the decision to treat the prototype as a supplement to lectures, diagrams, whiteboard explanation, and conventional practice materials.

\section{Environment Structure}

\textit{What should go here:} describe the overall layout of the environment and the logic of the learner's movement through it. Explain why the prototype was organised as a sequence of world-space teaching stations rather than as a floating menu or a purely abstract simulator.

\section{Mode Structure}

\textit{What should go here:} explain the roles of Learning, Practice, and Test. This section should show how the three modes embody progression from explicit guidance to more independent completion while preserving a common lesson structure.

\section{Interaction Model}

\textit{What should go here:} describe the main learner actions that matter pedagogically, such as decoding, register selection, datapath progression, operand or route choice, validation, and end-state reasoning. Make clear which interactions were made physical and which were left to UI support, and why.

\section{Representational and Pedagogical Abstractions}

\textit{What should go here:} explain where the prototype stays close to the textbook single-cycle datapath and where it deliberately departs from strict hardware realism. This is the place to discuss externalised hidden processes, staged progression, physical datapackets, and the decision to preserve reasoning structure rather than literal timing.

\section{Core Design Decisions and Trade-Offs}

\textit{What should go here:} gather the important design choices in one place. This should include scope narrowing, supported instruction classes, readability versus realism, stability versus feature expansion, and the trade-off between pedagogical clarity and simulation fidelity.
